\documentclass[12pt]{article}
\usepackage[utf8]{inputenc}
\usepackage[T1]{fontenc}
\usepackage{hyperref}
\title{Asmat, Central Language Resources for K-12 Curriculum Development}
\author{Compiled by Claude Sonnet 5.5 (claude-sonnet-5-5) for Andr\'as Kornai}
\date{}
\begin{document}
\maketitle

\section{Introduction}
Central Asmat (ISO 639-3 cns; Glottocode cent2117) is a Trans-New Guinea language of southern Papua, Indonesia, with roughly 7,000--10,000 speakers depending on the source. Searches found scholarly descriptions and a scripture portal, but no school materials, primers or curriculum documents. Little pedagogical material appears to exist online.

\section{Primary Educational Resources}
No primer, textbook or school programme for Central Asmat was found in the search results.

\section{Linguistic and Reference Materials}
\begin{itemize}
\item Overview page for Central Asmat at the New Guinea World project (University of Zurich): \url{https://newguineaworld.linguistik.uzh.ch/families/trans-new-guinea/asmat-muli-strait/asmat-kamrau-bay/asmat-kamoro/asmat/central-asmat}
\item C. L. Voorhoeve (1965), a study of the Flamingo Bay (Kawenak) dialect of Asmat, catalogued at Glottolog: \url{https://glottolog.org/resource/reference/id/137663}
\item C. L. Voorhoeve (1980), \emph{The Asmat Languages of Irian Jaya} (survey), catalogued at Glottolog: \url{https://glottolog.org/resource/reference/id/85998}
\item P. Drabbe (1959), Asmat dictionary, catalogued at Glottolog: \url{https://glottolog.org/resource/reference/id/475473}
\item ASJP word-list entry for Central Asmat: \url{https://asjp.clld.org/languages/ASMAT_CENTRAL}
\end{itemize}

\section{Religious and Community Resources}
\begin{itemize}
\item Scripture Earth listing for cns, which according to the search summary points to a Bible.com (YouVersion) text and Global Recordings Network audio: \url{https://scriptureearth.org/00eng.php?iso=cns}
\end{itemize}
No community organisations were found.

\section{Digital and Audio Resources}
The only audio indicated is the Global Recordings Network material reachable through the Scripture Earth page above. Nothing else was found.

\section{Curriculum Development Context}
The language is not known to be taught in schools, and no curriculum materials were found. A curriculum would have to be built from the scholarly descriptions above, which are older and cover mainly particular dialects (Flamingo Bay/Kawenak, Simai), together with new work with the community. Access to the printed works would need to be checked separately.

\end{document}